%% kapa: exam=1 round=29 year=2018 period=2 form=A subject=appraisal_law subject_ko=감정평가관계법규 seq=34 number=34
%% source: Q-Net 감정평가사 기출문제 (cst00309) artlSeq=5202808 file=extracted/1008/5202808/2018년 제29회 감정평가사 1차시험 기출문제/2018년 제29회 감정평가사 1차시험 기출문제/2교시 A형.pdf
%% conversion: kordoc 4.24.51 -> tools/gen1.py (auto); review=proofread-sonnet; flags=none
\begin{question}{34}
부동산등기법상 등기절차에 관한 설명으로 옳지 않은 것은?
\begin{choices}
  \item 대표자나 관리인이 있는 법인 아닌 사단이나 재단에 속하는 부동산의 등기에 관하여는 그 대표자나 관리인을 등기권리자 또는 등기의무자로 한다.
  \item 등기기록에 기록된 사항이 많아 취급하기에 불편하게 되는 등 합리적 사유로 등기기록을 옮겨 기록할 필요가 있는 경우에 등기관은 현재 효력이 있는 등기만을 새로운 등기기록에 옮겨 기록할 수 있다.
  \item 채권자는「민법」의 채권자대위권 규정에 따라 채무자를 대위하여 등기를 신청할 수 있다.
  \item 행정구역이 변경되었을 때에는 등기기록에 기록된 행정구역에 대하여 변경등기가 있는 것으로 본다.
  \item 등기는 법률에 다른 규정이 없는 한 당사자의 신청 또는 관공서의 촉탁에 따라 한다.
\end{choices}
\end{question}
\answer{1}
